约定
\[
  \binom{n}{k}=0\qquad(k<0\text{ 或 }k>n),
\]
因此下列求和可以直接遍历所有整数范围内有意义的项。

代码使用 Fermat 小定理求逆元，要求 \texttt{mod} 为质数且预处理上界小于
\texttt{mod}。
调用 \texttt{A}、\texttt{C} 时的第一个参数不能超过预处理上界。
常见模数 $10^9+7$、$998244353$ 可以直接用 \texttt{long long} 相乘；更大模数需要防止乘法溢出。

\subsubsection{基本恒等式}

\[
  \binom{n}{k}=\binom{n}{n-k},
  \qquad
  \binom{n}{k}=\binom{n-1}{k}+\binom{n-1}{k-1}.
\]

\[
  k\binom{n}{k}=n\binom{n-1}{k-1},
  \qquad
  (n-k)\binom{n}{k}=n\binom{n-1}{k}.
\]

\[
  \binom{n}{k}\binom{k}{m}
  =\binom{n}{m}\binom{n-m}{k-m}.
\]

\subsubsection{十二重计数法：球盒模型}

设有 $n$ 个球和 $k$ 个盒，$n,k\ge 1$。
三种限制分别为：任意放置、每盒至多一个球、每盒至少一个球。

第二类斯特林数 $S_2(n,k)$ 表示将 $n$ 个不同元素划分成恰好 $k$ 个非空、无标号集合的方案数。
边界与递推为
\[
  S_2(0,0)=1,
  \qquad
  S_2(n,0)=S_2(0,k)=0\quad(n,k>0),
\]
\[
  S_2(n,k)=S_2(n-1,k-1)+kS_2(n-1,k).
\]
给 $k$ 个集合重新编号，便得到将不同球放满不同盒的方案数
\[
  k!S_2(n,k)
  =\sum_{i=0}^{k}(-1)^i\binom{k}{i}(k-i)^n.
\]
第一类斯特林数计数的是具有指定环数的置换，本球盒模型使用第二类。

再记 $(k)_n=k!/(k-n)!$ 为下降阶乘，当 $n>k$ 时取 $0$；
记 $p_k(n)$ 为将 $n$ 表示成恰好 $k$ 个正整数之和的无序方案数。

\textbf{球不同，盒不同}
\[
\begin{aligned}
  \text{任意放置：}&\quad k^n,\\
  \text{每盒至多一个：}&\quad (k)_n,\\
  \text{每盒至少一个：}&\quad k!S_2(n,k).
\end{aligned}
\]

\textbf{球相同，盒不同}
\[
\begin{aligned}
  \text{任意放置：}&\quad \binom{n+k-1}{k-1},\\
  \text{每盒至多一个：}&\quad \binom{k}{n},\\
  \text{每盒至少一个：}&\quad \binom{n-1}{k-1}.
\end{aligned}
\]

\textbf{球不同，盒相同}
\[
\begin{aligned}
  \text{任意放置：}&\quad \sum_{i=0}^{k}S_2(n,i),\\
  \text{每盒至多一个：}&\quad
    \begin{cases}1,&n\le k,\\0,&n>k,\end{cases}\\
  \text{每盒至少一个：}&\quad S_2(n,k).
\end{aligned}
\]

\textbf{球相同，盒相同}
\[
\begin{aligned}
  \text{任意放置：}&\quad \sum_{i=0}^{k}p_i(n),\\
  \text{每盒至多一个：}&\quad
    \begin{cases}1,&n\le k,\\0,&n>k,\end{cases}\\
  \text{每盒至少一个：}&\quad p_k(n).
\end{aligned}
\]
其中不满足数量限制的其他式子也视为 $0$，且
\[
  p_k(n)=p_{k-1}(n-1)+p_k(n-k).
\]

\subsubsection{普通求和}

\[
  \sum_{k=0}^{n}\binom{n}{k}=2^n.
\]

\[
  \sum_{k=0}^{n}(-1)^k\binom{n}{k}
  =
  \begin{cases}
    1,&n=0,\\
    0,&n>0.
  \end{cases}
\]

当 $n>0$ 时，奇数项与偶数项各占一半：
\[
  \sum_{\substack{0\le k\le n\\k\text{ 为偶数}}}\binom{n}{k}
  =
  \sum_{\substack{0\le k\le n\\k\text{ 为奇数}}}\binom{n}{k}
  =2^{n-1}.
\]

当 $n\ge 2$ 时，常用带权和为
\[
  \sum_{k=0}^{n}k\binom{n}{k}=n2^{n-1},
\]
\[
  \sum_{k=0}^{n}k(k-1)\binom{n}{k}
  =n(n-1)2^{n-2},
\]
\[
  \sum_{k=0}^{n}k^2\binom{n}{k}
  =n(n+1)2^{n-2}.
\]

\subsubsection{固定下指标连续求和}

冰球杆恒等式为
\[
  \sum_{i=k}^{n}\binom{i}{k}=\binom{n+1}{k+1}.
\]
因此区间和可以写成
\[
  \sum_{i=l}^{r}\binom{i}{k}
  =\binom{r+1}{k+1}-\binom{l}{k+1}.
\]

\subsubsection{范德蒙德恒等式}

\[
  \sum_k\binom{n}{k}\binom{m}{r-k}
  =\binom{n+m}{r}.
\]

常用特例和带权形式为
\[
  \sum_{k=0}^{n}\binom{n}{k}^2=\binom{2n}{n},
\]
\[
  \sum_k k\binom{n}{k}\binom{m}{r-k}
  =n\binom{n+m-1}{r-1}.
\]

\subsubsection{固定上指标连续求和}

定义组合数前缀和
\[
  S(n,k)=\sum_{i=0}^{k}\binom{n}{i}.
\]
由 Pascal 恒等式可得递推
\[
  S(n,k)=S(n-1,k)+S(n-1,k-1)
  =2S(n-1,k)-\binom{n-1}{k}.
\]
任意连续区间可以用前缀和相减：
\[
  \sum_{i=l}^{r}\binom{n}{i}=S(n,r)-S(n,l-1).
\]

\subsubsection{模 $2$ 的连续和}

当 $n\ge 1$ 且 $0\le k\le n$ 时，由前缀和递推立即得到
\[
  \sum_{i=0}^{k}\binom{n}{i}
  \equiv\binom{n-1}{k}\pmod 2.
\]
由 Lucas 定理，组合数的奇偶性满足
\[
  \binom{n}{k}\equiv1\pmod2
  \iff (k\mathbin{\&}n)=k.
\]
所以
\[
  \sum_{i=0}^{k}\binom{n}{i}\equiv1\pmod2
  \iff (k\mathbin{\&}(n-1))=k.
\]
当 $0\le l\le r\le n$ 时，区间版本为
\[
  \sum_{i=l}^{r}\binom{n}{i}
  \equiv
  \binom{n-1}{r}+\binom{n-1}{l-1}
  \pmod2.
\]

\subsubsection{杂项}

\textbf{循环引理}

设整数序列 $a_1,a_2,\ldots,a_n$ 满足
\[
  a_i\le 1,
  \qquad
  \sum_{i=1}^{n}a_i=k>0.
\]
把每个位置都视为一个循环起点，并按起点位置计数。
恰有 $k$ 个循环移位满足它的每个非空前缀和都严格大于 $0$。
特别地，当序列总和为 $1$ 时，恰有一个满足条件的循环移位；
它从原序列前缀和最后一次取得全局最小值的位置之后开始。

\textbf{Ballot 定理}

候选人 A 最终得到 $p$ 票，候选人 B 最终得到 $q$ 票，且 $p>q$。
把投给 A、B 的票分别记为 $+1,-1$，则 A 在计票过程中始终严格领先的序列数为
\[
  \frac{p-q}{p+q}\binom{p+q}{q}.
\]
若只要求 A 从不落后，允许中途平票，则当 $p\ge q$ 时方案数为
\[
  \binom{p+q}{q}-\binom{p+q}{q-1}
  =\frac{p-q+1}{p+1}\binom{p+q}{q}.
\]

\textbf{Catalan 特例}

当 $p=q=n$ 时，从不落后的序列数为
\[
  C_n=\frac{1}{n+1}\binom{2n}{n}.
\]
它也计数长度为 $2n$ 的合法括号序列，以及从 $(0,0)$ 走到 $(n,n)$ 且不越过对角线的格路。

\textbf{反射原理}

对于从 $(0,0)$ 走到 $(p,q)$ 的向右、向上格路，设 $p\ge q$。
将第一次到达 $y=x+1$ 之前的路段关于 $y=x$ 反射，
可将越过 $y=x$ 的坏路与 $\binom{p+q}{q-1}$ 条格路一一对应。
因此不越过 $y=x$ 的格路数为
\[
  \binom{p+q}{q}-\binom{p+q}{q-1}
  =\frac{p-q+1}{p+1}\binom{p+q}{q}.
\]
一般地，遇到“路径不得穿过一条直线”时，可按首次碰到边界的位置反射前缀，用总路径数减去反射后的坏路径数。

\textbf{LGV 引理}

设 $G$ 是有向无环图，路径的权值为路上边权之积。
记 $e(A_i,B_j)$ 为从起点 $A_i$ 到终点 $B_j$ 的所有路径权值之和。
令 $\mathcal P_\sigma$ 为所有满足 $P_i:A_i\to B_{\sigma(i)}$ 且两两点不交的路径组，则
\[
  \begin{aligned}
    D
    &=\sum_{\sigma\in S_r}\operatorname{sgn}(\sigma)
      \sum_{(P_1,\ldots,P_r)\in\mathcal P_\sigma}
      \prod_{i=1}^{r}w(P_i)\\
    &=
    \begin{vmatrix}
      e(A_1,B_1) & e(A_1,B_2) & \cdots & e(A_1,B_r)\\
      e(A_2,B_1) & e(A_2,B_2) & \cdots & e(A_2,B_r)\\
      \vdots & \vdots & \ddots & \vdots\\
      e(A_r,B_1) & e(A_r,B_2) & \cdots & e(A_r,B_r)
    \end{vmatrix}.
  \end{aligned}
\]
无权时，每组路径的权值都为 $1$。

\textbf{MacMahon 盒装平面分拆公式}

平面分拆可写成一个行、列均弱递减的非负整数数组
$\pi=(\pi_{ij})$。
能装入 $a\times b\times c$ 盒子，即数组不超过 $a$ 行、$b$ 列，且
$0\le\pi_{ij}\le c$ 的平面分拆数为
\[
  \prod_{i=1}^{a}\prod_{j=1}^{b}\prod_{k=1}^{c}
  \frac{i+j+k-1}{i+j+k-2}
  =\prod_{i=1}^{a}\prod_{j=1}^{b}
  \frac{i+j+c-1}{i+j-1}.
\]
它也等于边长依次为 $a,b,c,a,b,c$ 的六边形的菱形铺砌数。
上式整体一定为整数；模意义下只有当分母可逆时才能逐项求逆，
否则需先约分或统计质因子次数。
